\begin{figure}[H]
\centering
\resizebox{\textwidth}{!}{%
\begin{tikzpicture}[
    node distance=1.2cm and 1.8cm,
    >=Stealth, thick, font=\sffamily
]

\tikzset{
    io_node/.style={font=\sffamily\bfseries},
    dot/.style={circle, fill=black, inner sep=0pt, minimum size=4pt}
}

% ---------------------------------------------------------
% Bloco de Processamento de Entrada (esquerda)
% ---------------------------------------------------------
\node[block, minimum width=4.0cm, minimum height=4.0cm, text width=3.5cm] (IPB) at (0, 0) {%
    \textbf{Input}\\
    \textbf{Processing}\\
    \textbf{Block}\\[3pt]
    {\footnotesize CORDIC + Phase}\\
    {\footnotesize Unwrap}\\
    {\footnotesize (40 ciclos)}
};

% ---------------------------------------------------------
% Entradas (extremo esquerdo)
% ---------------------------------------------------------
\node[io_node] (Iin) at (-4.5, 1.0) {$i_{in}$};
\node[io_node] (Qin) at (-4.5, -0.5) {$q_{in}$};
\node[io_node] (Vin) at (-4.5, -2.0) {input\_valid};

\draw[data arrow] (Iin) -- ($(IPB.west) + (0, 1.0)$);
\draw[data arrow] (Qin) -- ($(IPB.west) + (0, -0.2)$);
\draw[ctrl arrow] (Vin) -- ($(IPB.west) + (0, -1.4)$);

% ---------------------------------------------------------
% Sa\'idas do IPB
% ---------------------------------------------------------
% envelope_s (topo direito do IPB)
\coordinate (envOut) at ($(IPB.east) + (0, 1.2)$);
% theta_norm (meio direito do IPB)
\coordinate (thetaNormOut) at ($(IPB.east) + (0, 0.0)$);
% phase_wrapped (baixo direito do IPB)
\coordinate (phaseWOut) at ($(IPB.east) + (0, -1.2)$);
% buffers_enable_raw (base do IPB)
\coordinate (enRawOut) at ($(IPB.south) + (0.5, 0)$);

% ---------------------------------------------------------
% LUTs trigonometricas (secao central)
% ---------------------------------------------------------
\node[block, minimum width=2.8cm, minimum height=1.5cm, text width=2.4cm] (LUT_C) at (5.5, 1.5) {%
    \textbf{LUT\_Trig}\\
    {\footnotesize (cosseno)}\\
    {\footnotesize $\cos(\Delta\theta)$}
};

\node[block, minimum width=2.8cm, minimum height=1.5cm, text width=2.4cm] (LUT_S) at (5.5, -1.5) {%
    \textbf{LUT\_Trig}\\
    {\footnotesize (seno)}\\
    {\footnotesize $\sin(\Delta\theta)$}
};

% theta_norm bifurca para ambas LUTs
\draw[data arrow] (thetaNormOut) -- ++(1.5, 0) coordinate (thetaFork);
\node[dot] at (thetaFork) {};
\draw[data arrow] (thetaFork) |- (LUT_C.west)
    node[signal, pos=0.3, above] {\scriptsize $\theta_{norm}$};
\draw[data arrow] (thetaFork) |- (LUT_S.west);

% ---------------------------------------------------------
% Bloco de Atraso de Alinhamento (envelope + theta)
% ---------------------------------------------------------
\node[small block, minimum width=2.8cm, minimum height=1.8cm, text width=2.4cm] (Delay) at (5.5, 4.5) {%
    \textbf{Atraso de}\\
    \textbf{Alinhamento}\\
    {\scriptsize 4 ciclos (LUT)}
};

% envelope_s -> Atraso
\draw[data arrow] (envOut) -- ++(1.5, 0) |- ($(Delay.west) + (0, 0.3)$)
    node[signal, pos=0.2, above] {\scriptsize envelope\_s};

% phase_wrapped -> Atraso (caminho theta)
\draw[data arrow] (phaseWOut) -- ++(0.8, 0) |- ($(Delay.west) + (0, -0.3)$)
    node[signal, pos=0.15, below] {\scriptsize phase\_wrapped};

% ---------------------------------------------------------
% Atraso do Enable
% ---------------------------------------------------------
\node[small block, minimum width=2.8cm, minimum height=0.8cm, text width=2.4cm] (EnDelay) at (5.5, -4.0) {%
    \textbf{Atraso Enable}\\
    {\scriptsize 4 ciclos}
};
\draw[ctrl arrow] (enRawOut) |- (EnDelay.west)
    node[signal, pos=0.3, below] {\scriptsize en\_raw};

% ---------------------------------------------------------
% Registradores de Deslocamento (secao direita)
% ---------------------------------------------------------
\node[block, minimum width=3.2cm, minimum height=1.3cm, text width=2.8cm] (SR_E) at (11.5, 4.5) {%
    \textbf{Shift\_Reg}\\
    {\footnotesize (envelope, $M\!=\!5$)}
};

\node[block, minimum width=3.2cm, minimum height=1.3cm, text width=2.8cm] (SR_C) at (11.5, 1.5) {%
    \textbf{Shift\_Reg}\\
    {\footnotesize (cosseno, $M\!=\!5$)}
};

\node[block, minimum width=3.2cm, minimum height=1.3cm, text width=2.8cm] (SR_S) at (11.5, -1.5) {%
    \textbf{Shift\_Reg}\\
    {\footnotesize (seno, $M\!=\!5$)}
};

% Atraso -> Shift_Reg (envelope)
\draw[data arrow] (Delay.east) -- ++(0.5, 0) |- (SR_E.west)
    node[signal, pos=0.4, above] {\scriptsize env\_aligned};

% LUT_C -> Shift_Reg (cos)
\draw[data arrow] (LUT_C.east) -- (SR_C.west)
    node[signal, midway, above] {\scriptsize cos\_dp};

% LUT_S -> Shift_Reg (sin)
\draw[data arrow] (LUT_S.east) -- (SR_S.west)
    node[signal, midway, above] {\scriptsize sin\_dp};

% Distribuicao do enable
\coordinate (enAligned) at ($(EnDelay.east) + (1.5, 0)$);
\draw[ctrl arrow] (EnDelay.east) -- (enAligned);
\node[dot] at (enAligned) {};
\draw[ctrl arrow] (enAligned) |- ($(SR_E.south) + (-0.5, 0)$)
    node[signal, pos=0.1, right] {\scriptsize wr\_en};
\draw[ctrl arrow] (enAligned) |- ($(SR_C.south) + (-0.5, 0)$);
\draw[ctrl arrow] (enAligned) |- ($(SR_S.south) + (-0.5, 0)$);

% ---------------------------------------------------------
% Saidas (extremo direito)
% ---------------------------------------------------------
\node[io_node, right=1.5cm of SR_E] (OutE) {buf\_env\_out$[0\!:\!5]$};
\node[io_node, right=1.5cm of SR_C] (OutC) {buf\_cos\_out$[0\!:\!5]$};
\node[io_node, right=1.5cm of SR_S] (OutS) {buf\_sin\_out$[0\!:\!5]$};

\draw[data arrow] (SR_E.east) -- (OutE) node[bus, midway, below] {$/6$};
\draw[data arrow] (SR_C.east) -- (OutC) node[bus, midway, below] {$/6$};
\draw[data arrow] (SR_S.east) -- (OutS) node[bus, midway, below] {$/6$};

% ---------------------------------------------------------
% theta_unwrapped_out (CORRIGIDO: acima de buf_env, sem sobreposicao)
% ---------------------------------------------------------
% Posicionado bem acima da linha do SR_E para evitar colisao visual
\node[io_node] (OutTheta) at ($(SR_E.north) + (3.5, 1.8)$) {theta\_unwrapped\_out};

% Fio: do lado leste do Delay, sobe ate a saida de theta
\draw[data arrow] ($(Delay.east) + (0, -0.3)$) -- ++(2.0, 0) coordinate (ThetaMid);
\draw[data arrow] (ThetaMid) -- ($(ThetaMid |- OutTheta)$) -- (OutTheta);
\node[signal] at ($(ThetaMid) + (0, 0.3)$) {\scriptsize $\theta_{aligned}$};

% ---------------------------------------------------------
% stage_valid
% ---------------------------------------------------------
\node[io_node] (OutValid) at ($(enAligned) + (4.5, 0)$) {stage\_valid};
\draw[ctrl arrow] (enAligned) -- ++(0.5, 0) |- (OutValid);

% ---------------------------------------------------------
% Borda da entidade
% ---------------------------------------------------------
\node[entity border, fit=(Iin)(Vin)(OutE)(OutS)(OutTheta)(OutValid)(Delay)(EnDelay), inner sep=15pt] {};

\end{tikzpicture}%
}

\vspace{3mm}
\noindent\textit{Lat\^{e}ncia das LUTs $= 4$ ciclos. O atraso de alinhamento garante sincronismo do envelope e de $\theta$ com as sa\'{i}das das LUTs. Todos os registradores de deslocamento avancam simultaneamente com \textsf{wr\_en}.}

\caption{Top\_NeuralNet\_S1 -- Est\'{a}gio 1: convers\~{a}o IQ$\,\rightarrow\,$polar, decomposi\c{c}\~{a}o trigonom\'{e}trica e buffers deslizantes.}
\label{fig:top_neuralnet_s1}
\end{figure}
